\section{Free pseudo-product fundamental graded Lie algebras}\label{section8}

An FGLA $\mathfrak m=\bigoplus\limits_{p<0}\mathfrak g_p$ equipped with
nonzero subspaces $\mathfrak e$, $\mathfrak f$ of $\mathfrak g_{-1}$ is called
a pseudo-product FGLA if the following conditions hold:
\begin{enumerate}\itemsep=0pt
\renewcommand{\labelenumi}{$(\roman{enumi})$}
\item $\mathfrak g_{-1}=\mathfrak e\oplus\mathfrak f$;
\item $[\mathfrak e,\mathfrak e]=[\mathfrak f,\mathfrak f]=\{0\}$.
\end{enumerate}
The pair $(\mathfrak e,\mathfrak f)$ is called the pseudo-product structure of
the pseudo-product FGLA
$\mathfrak m=\bigoplus\limits_{p<0}\mathfrak g_p$.
We will also denote by the triplet $(\mathfrak m;\mathfrak e,\mathfrak f)$
the pseudo-product FGLA $\mathfrak m=\bigoplus\limits_{p<0}\mathfrak g_p$
with pseudo-product structure $(\mathfrak e,\mathfrak f)$.
Let $\mathfrak m=\bigoplus\limits_{p<0}\mathfrak g_p$ (resp.
$\mathfrak m'=\bigoplus\limits_{p<0}\mathfrak g'_p$) be a
pseudo-product FGLA with pseudo-product structure $(\mathfrak e,\mathfrak f)$
(resp. $(\mathfrak e',\mathfrak f')$).
We say that two pseudo-product FGLAs $(\mathfrak m;\mathfrak e,\mathfrak f)$
and $(\mathfrak m';\mathfrak e',\mathfrak f')$
are isomorphic if there exists a GLA isomorphism
$\varphi$ of $\mathfrak m$ onto $\mathfrak m'$ such that
$\varphi(\mathfrak e)=\mathfrak e'$ and $\varphi(\mathfrak f)=\mathfrak f'$.
\begin{proposition}\label{prop8.1} Let $\mathfrak m=\bigoplus\limits_{p<0}\mathfrak g_p$ be a
pseudo-product FGLA of the $\mu$-th kind with pseudo-product
structure $(\mathfrak e,\mathfrak f)$.
If $\mathfrak m$ is a free FGLA of type $(n,\mu)$,
then $n=2$.
\end{proposition}

\begin{proof} Let $(e_1,\dots,e_m)$ (resp.\ $(f_1,\dots,f_l)$) be a basis of
$\mathfrak e$ (resp.~$\mathfrak f$).
Since $[\mathfrak e,\mathfrak f]=\mathfrak g_{-2}$,
the space~$\mathfrak g_{-2}$ is generated by
$\{[e_i,f_j]:i=1,\dots,m,j=1,\dots,l\}$ as a vector space, so
$\dim \mathfrak g_{-2}\leqq ml$. On the other hand,
since $\mathfrak m$ is a free FGLA,
\[
\dim \mathfrak g_{-2}=\dim b(\mathfrak g_{-1},\mu)_{-2}=
\dim \Lambda^2(\mathfrak g_{-1})=\frac{(m+l)(m+l-1)}{2},
\]
so $ml\geqq \frac{(m+l)(m+l-1)}{2}$.
From this fact it follows that $m=l=1$.
\end{proof}

Let $\mathfrak m=\bigoplus\limits_{p<0}\mathfrak g_p$ be a pseudo-product FGLA
of the $\mu$-th kind with
pseudo-product structure $(\mathfrak e,\mathfrak f)$, where $\mu\geqq2$.
$\mathfrak m$ is called a free pseudo-product FGLA
of type $(m,n,\mu)$ if the following conditions hold:
\begin{enumerate}\itemsep=0pt
\renewcommand{\labelenumi}{$(\roman{enumi})$}
\item $\dim\mathfrak e=m$ and $\dim\mathfrak f=n$;
\item Let $\mathfrak m'=\bigoplus\limits_{p<0}\mathfrak g'_p$ be a pseudo-product FGLA of the $\mu$-th kind with
pseudo-product structure $(\mathfrak e',\mathfrak f')$
and let $\varphi$ be a surjective linear mapping of $\mathfrak g_{-1}$ onto
$\mathfrak g'_{-1}$ such that
$\varphi(\mathfrak e)\subset\mathfrak e'$ and
$\varphi(\mathfrak f)\subset\mathfrak f'$.
Then $\varphi$ can be extended uniquely to a GLA
epimorphism of $\mathfrak m$ onto $\mathfrak m'$.
\end{enumerate}
\begin{proposition}\label{prop8.2} Let $m$, $n$ and $\mu$ be positive integers such that $\mu\geqq2$.
\begin{enumerate}\itemsep=0pt
\item There exists a unique free pseudo-product FGLA
of type
$(m,n,\mu)$ up to isomorphism.
\item Let $\mathfrak m=\bigoplus\limits_{p<0}\mathfrak g_p$ be a free
pseudo-product FGLA of type $(m,n,\mu)$
with pseudo-product structure $(\mathfrak e,\mathfrak f)$.
We denote by $\Der(\mathfrak m;\mathfrak e,\mathfrak f)_0$ the Lie algebra of
all the derivations of $\mathfrak m$ preserving the gradation of $\mathfrak m$, $\mathfrak e$ and $\mathfrak f$.
Then the mapping $\Phi:\Der(\mathfrak m;\mathfrak e,\mathfrak f)_0\ni D\mapsto (D|\mathfrak e,D|\mathfrak f)\in\gl(\mathfrak e)\times\gl(\mathfrak f)$ is a Lie algebra isomorphism.
\end{enumerate}
\end{proposition}

\begin{proof} (1)
The uniqueness of a free pseudo-product FGLA of type $(m,n,\mu)$ follows from the def\/inition.
Let $V$ be an $(m+n)$-dimensional vector space and let
$\mathfrak e$, $\mathfrak f$ be subspaces of $V$
such that $V=\mathfrak e\oplus \mathfrak f$,
$\dim \mathfrak e=m$ and $\dim \mathfrak f=n$.
Let $\mathfrak a=\bigoplus\limits_{p<0}\mathfrak a_p$ be the graded ideal
of $b(V,\mu)$ generated by
$[\mathfrak e,\mathfrak e]+[\mathfrak f,\mathfrak f]$.
We set $\mathfrak m=b(V,\mu)/\mathfrak a$,
$\mathfrak g_p=b(V,\mu)_p/\mathfrak a_p$.
Clearly $\mathfrak m=\bigoplus\limits_{p<0}\mathfrak g_p$ is
a~pseudo-product FGLA.
We show that the factor algebra $\mathfrak m$
is a free pseudo-product FGLA of type $(m,n,\mu)$.
First we prove that $\mathfrak m$ is of the $\mu$-th kind.
Let $\mathfrak n=\bigoplus\limits_{p<0}\mathfrak g''_p$ be
a free FGLA of type $(2,\mu)$ and let~$\mathfrak e''$ and~$\mathfrak f''$
be one-dimensional subspaces of $\mathfrak g''_{-1}$ such that
$\mathfrak g''_{-1}=\mathfrak e''\oplus\mathfrak f''$.
Let~$\varphi_1$ be an injective linear mapping of $\mathfrak g''_{-1}$
into $V$ such that
$\varphi_1(\mathfrak e'')\subset \mathfrak e$ and
$\varphi_1(\mathfrak f'')\subset \mathfrak f$.
Let~$\varphi_2$ be a linear mapping of~$V$
into $\mathfrak g''_{-1}$ such that
$\varphi_2\circ\varphi_1=1_{\mathfrak g''_{-1}}$,
$\varphi_2(\mathfrak e)=\mathfrak e''$ and
$\varphi_2(\mathfrak f)=\mathfrak f''$.
There exists
a~homomorphism~$L(\varphi_1)$ (resp.~$L(\varphi_2)$) of
$\mathfrak n$ (resp.~$b(V,\mu)$) into
$b(V,\mu)$ (resp.~$\mathfrak n$)
such that
$L(\varphi_1)|\mathfrak g''_{-1}=\varphi_1$
(resp.\ $L(\varphi_2)|V=\varphi_2$).
Since $L(\varphi_2)([\mathfrak e,\mathfrak e]+[\mathfrak f,\mathfrak f])=\{0\}$,
$L(\varphi_2)$ induces a~homomorphism $\hat{L}(\varphi_2)$ of
$\mathfrak m$ into $\mathfrak n$ such that
$L(\varphi_2)=\hat{L}(\varphi_2)\circ \pi$,
where $\pi$ is the natural projection of $b(V,\mu)$ onto
$\mathfrak m$. Since
\[
1_{\mathfrak n}=L(\varphi_2)\circ L(\varphi_1)
=\hat{L}(\varphi_2)\circ \pi\circ L(\varphi_1),
\]
$\pi\circ L(\varphi_1)$ is a monomorphism of $\mathfrak n$ into $\mathfrak m$, so $\mathfrak g_{-\mu}\ne\{0\}$.
Thus $\mathfrak m$ is of the $\mu$-th kind.
Let $\mathfrak m'=\bigoplus\limits_{p<0}\mathfrak g'_p$ be a pseudo-product
FGLA of the $\mu$-th kind
with pseudo-product structure $(\mathfrak e',\mathfrak f')$
and let $\phi$ be a surjective linear mapping of $b(V,\mu)_{-1}$ onto
$\mathfrak g'_{-1}$ such that
$\phi(\mathfrak e)\subset\mathfrak e'$ and
$\phi(\mathfrak f)\subset\mathfrak f'$.
By the def\/inition of a free FGLA,
there exists a GLA epimorphism $L(\phi)$ of
$b(V,\mu)$ onto $\mathfrak m'$ such that
$L(\phi)|b(V,\mu)_{-1}=\phi$.
Since $L(\phi)([\mathfrak e,\mathfrak e]+[\mathfrak f,\mathfrak f])\subset [\mathfrak e',\mathfrak e']+[\mathfrak f',\mathfrak f']=\{0\}$,
we see that $L(\phi)(\mathfrak a)=\{0\}$, so
the epimorphism $L(\phi)$ induces a GLA epimorphism
$\hat{L}(\phi)$ of $\mathfrak m$ onto $\mathfrak m'$ such that
$\hat{L}(\phi)|\mathfrak g_{-1}=\phi$.

(2)~We may prove the fact that the mapping $\Phi$ is surjective.
Let $\phi$ be an endomorphism of~$\mathfrak g_{-1}$ such that
$\phi(\mathfrak e)\subset \mathfrak e$ and
$\phi(\mathfrak f)\subset \mathfrak f$.
By Proposition~\ref{prop2.1}~(2), there exists a $D\in\Der(b(V,\mu))_0$
such that $D|b(V,\mu)_{-1}=\phi$.
Since $D([\mathfrak e,\mathfrak e]+[\mathfrak f,\mathfrak f])
\subset [\mathfrak e,\mathfrak e]+[\mathfrak f,\mathfrak f]$,
$D$ induces a derivation~$\hat{D}$ of~$\mathfrak m$ such that
$\hat{D}|\mathfrak g_{-1}=\phi$.
\end{proof}

\begin{remark}\label{rem8.1}
Let $m$, $n$, $m'$, $n'$ and $\mu$ be positive integers with $\mu\geqq2$, and let
$\mathfrak m=\bigoplus\limits_{p<0}\mathfrak g_p$
(resp.\ $\mathfrak m'=\bigoplus\limits_{p<0}\mathfrak g'_p$) be a
free pseudo-product FGLA of type $(m,n,\mu)$ $($resp.~$(m',n',\mu))$
with pseudo-product structure $(\mathfrak e,\mathfrak f)$
(resp. $(\mathfrak e',\mathfrak f'))$.
Furthermore
let $\varphi$ be a linear mapping of $\mathfrak g_{-1}$ into
$\mathfrak g'_{-1}$ such that $\varphi(\mathfrak e)\subset \mathfrak e'$ and
$\varphi(\mathfrak f)\subset \mathfrak f'$.
\begin{enumerate}\itemsep=0pt
\item
From the proof of Proposition~\ref{prop8.2}, there exists a unique GLA
homomorphism $\hat{L}(\varphi)$ of $\mathfrak m$ into $\mathfrak m'$
such that $\hat{L}(\varphi)|\mathfrak g_{-1}=\varphi$.
If $\varphi$ is injective, then $\hat{L}(\varphi)$ is a monomorphism.
\item
Assume that $m=n=1$ and $\varphi$ is injective.
Then $\hat{L}(\varphi)(\mathfrak m)$ is a graded subalgebra of $\mathfrak m'$
isomorphic to a free FGLA of type $(2,\mu)$.
From this result,
the subalgebra of $\mathfrak m'$ generated by a nonzero element $X$ of
$\mathfrak e'$ and a nonzero element $Y$ of $\mathfrak f'$
is a free FGLA of type $(2,\mu)$.
\end{enumerate}
\end{remark}

Let $\mathfrak m=\bigoplus\limits_{p<0}\mathfrak g_p$ be a pseudo-product FGLA
of the $\mu$-th kind with pseudo-product structure~$(\mathfrak e,\mathfrak f)$.
We denote by $\mathfrak g_0$ the Lie algebra of
all the derivations of $\mathfrak m$ preserving the gradation of~$\mathfrak m$,~$\mathfrak e$ and~$\mathfrak f$:
\[
\mathfrak g_0=\{ D\in\Der(\mathfrak g)_0:D(\mathfrak e)\subset \mathfrak e,
D(\mathfrak f)\subset \mathfrak f \}.
\]
The prolongation
$\gla g$ of $(\mathfrak m,\mathfrak g_0)$ is called the prolongation of
$(\mathfrak m;\mathfrak e,\mathfrak f)$.

A transitive GLA $\gla g$ is called a pseudo-product GLA
if there are given nonzero subspaces $\mathfrak e$ and $\mathfrak f$ of $\mathfrak g_{-1}$ satisfying the following conditions:
\begin{enumerate}\itemsep=0pt
\renewcommand{\labelenumi}{$(\roman{enumi})$}
\item The negative part $\mathfrak g_{-}$ is
a pseudo-product
FGLA with pseudo-product structure
$(\mathfrak e,\mathfrak f)$;
\item $[\mathfrak g_0,\mathfrak e]\subset \mathfrak e$
and $[\mathfrak g_0,\mathfrak f]\subset \mathfrak f$.
\end{enumerate}
The pair $(\mathfrak e,\mathfrak f)$ is called the pseudo-product structure of
 the pseudo-product GLA
$\gla g$.
If the $\mathfrak g_0$-modules $\mathfrak e$ and $\mathfrak f$ are
irreducible,
then the pseudo-product GLA $\gla g$ is said to be
of irreducible type.

The following lemma is due to N.~Tanaka (cf.~\cite{tan89:01}).
Here we give a proof for the convenience of the readers.
\begin{lemma}\label{lemma8.1} Let $\gla g$ be a pseudo-product GLA of
depth $\mu$ with pseudo-product structure~$(\mathfrak e,\mathfrak f)$.
\begin{enumerate}\itemsep=0pt
\item If $\mathfrak g_-$ is non-degenerate,
then $\mathfrak g$ is finite-dimensional.
\item If $\gla g$ is of irreducible type and $\mu\geqq2$, then
$\mathfrak g$ is finite-dimensional.
\end{enumerate}
\end{lemma}

\begin{proof} (1) The proof is essentially due to the proof of
\cite[Corollary 3 to Theorem 11.1]{tan70:1}.
For $p\in \mathbb Z$, we set
$\mathfrak h_p=\{X\in\mathfrak g_p:[X,\mathfrak g_{\leqq -2}]=\{0\}\}$.
We def\/ine $I\in\mathfrak{gl}(\mathfrak g_{-1})$ as follows:
$I(x)=-\sqrt{-1}x$ for $x\in\mathfrak e$,
$I(x)=\sqrt{-1}x$ for $x\in\mathfrak f$.
Then $I^2=-1$,
$I([a,x])=[a,I(x)]$ and $[I(x),I(y)]=[x,y]$ for $a\in\mathfrak g_0$ and
$x,y\in\mathfrak g_{-1}$. We put $\langle x,y\rangle=[I(x),y]$ for
$x,y\in\mathfrak g_{-1}$. Then $\langle x,y\rangle=\langle y,x\rangle$,
and for $x\in\mathfrak g_{-1}$, $\langle x,\mathfrak g_{-1}\rangle=\{0\}$
implies $x=0$, since $\mathfrak g_-$ is non-degenerate. Also
$\langle[a,x],y\rangle+\langle x,[a,y]\rangle=0$ and $[[b,x],y]=[[b,y],x]$
for $a\in\mathfrak h_0$, $b\in\mathfrak h_1$ and $x,y\in\mathfrak g_{-1}$.
Then, for $b\in\mathfrak h_1$,
$x,y,z\in\mathfrak g_{-1}$, we have $\langle[[b,x],y],z\rangle=
-\langle y,[[b,x],z]\rangle=-\langle y,[[b,z],x]\rangle=
\langle [[b,z],y],x\rangle=\langle[[b,y],z],x\rangle
=-\langle z,[[b,y],x]\rangle=-\langle[[b,x],y],z\rangle$, so
$\langle[[b,x],y],z\rangle=0$. By transitivity of $\mathfrak g$,
$\mathfrak h_1=\{0\}$. Therefore by \cite[Corollary 1 to Theorem 11.1]{tan70:1}, $\mathfrak g$ is f\/inite-dimensional.

(2) We may assume that $\mathfrak g_1\ne\{0\}$.
By \cite[Lemma 2.4]{yat88:0}, the
$\mathfrak g_0$-modules $\mathfrak e$, $\mathfrak f$ are not isomorphic to each
other. We put
$\mathfrak d=\{X\in\mathfrak g_{-1}:[X,\mathfrak g_{-1}]=\{0\}\}$;
then $\mathfrak d$ is a $\mathfrak g_0$-submodule of $\mathfrak g_{-1}$.
Hence $\mathfrak d=\{0\}$, $\mathfrak d=\mathfrak e$, $\mathfrak d=\mathfrak f$ or
$\mathfrak d=\mathfrak g_{-1}$. If $\mathfrak d\ne\{0\}$, then $\mathfrak g_{-2}=[\mathfrak e,\mathfrak f]=\{0\}$, which is a contradiction. Thus $\mathfrak g_-$ is non-degenerate.
By (1), $\mathfrak g$ is f\/inite-dimensional.
\end{proof}

The prolongation of a pseudo-product FGLA becomes
a pseudo-product GLA.
By Proposition~\ref{prop8.2}~(2), the prolongation of a free pseudo-product
FGLA is a pseudo-product GLA of irreducible type.
By Lemma \ref{lemma8.1}~(2), the prolongation of a free pseudo-product FGLA is
f\/inite-dimensional.

\begin{proposition}\label{prop8.3} Let $\mathfrak m=\bigoplus\limits_{p<0}\mathfrak g_p$ be a
free pseudo-product FGLA of type $(m,n,\mu)$
with pseudo-product structure $(\mathfrak e,\mathfrak f)$
and let
$\gla g$ be the prolongation of $(\mathfrak m;\mathfrak e,\mathfrak f)$.
\begin{enumerate}\itemsep=0pt
\item
$\mathfrak g_0$ is isomorphic to $\gl(\mathfrak e)\oplus\gl(\mathfrak f)$ as
a Lie algebra.
\item $\mathfrak g_{-2}$ is isomorphic to $\mathfrak e\otimes \mathfrak f$
as a $\mathfrak g_0$-module. In particular,
$\dim\mathfrak g_{-2}=mn$.
\item $\mathfrak g_{-3}$ is isomorphic to
$S^2(\mathfrak e)\otimes
\mathfrak f\oplus S^2(\mathfrak f)\otimes \mathfrak e$
as a $\mathfrak g_0$-module.
In particular,
$\dim\mathfrak g_{-3}=\frac{mn(m+n+2)}2$.
\end{enumerate}
\end{proposition}

\begin{proof} (1) This follows from Proposition \ref{prop8.2} (2).
\par
(2) Let $\mathfrak a=\bigoplus\limits_{p<0}\mathfrak a_p$ be the graded ideal
of $b(\mathfrak g_{-1},\mu)$ generated by
$[\mathfrak e,\mathfrak e]+[\mathfrak f,\mathfrak f]$.
By the construction of $b(\mathfrak g_{-1},\mu)_{-2}$,
$\mathfrak a_{-2}$ is isomorphic to $\Lambda^2(\mathfrak e)\oplus
\Lambda^2(\mathfrak f)$, so
$\mathfrak g_{-2}=b(\mathfrak g_{-1},\mu)_{-2}/\mathfrak a_{-2}$
is isomorphic to $\mathfrak e\otimes \mathfrak f$.

(3) By the construction of $b(\mathfrak g_{-1},\mu)_{-3}$,
$b(\mathfrak g_{-1},\mu)_{-3}$ is isomorphic to
\[
(\mathfrak e\oplus \mathfrak f)\otimes \Lambda^2(\mathfrak e\oplus \mathfrak f)/\Lambda^3(\mathfrak e\oplus \mathfrak f)\cong
(\mathfrak e\otimes \mathfrak e\otimes \mathfrak f)
\oplus
(\mathfrak e\otimes \mathfrak f\otimes \mathfrak f).
\]
Moreover,
$\mathfrak a_{-3}$ is isomorphic to
\[
(\mathfrak e\oplus \mathfrak f)\otimes \Lambda^2(\mathfrak e)
\oplus(\mathfrak e\oplus \mathfrak f)\otimes \Lambda^2(\mathfrak f)/
\Lambda^3(\mathfrak e\oplus \mathfrak f)
\cong \mathfrak e\otimes\Lambda^2(\mathfrak e)\oplus
\mathfrak f\otimes\Lambda^2(\mathfrak f).
\]
Hence
$\mathfrak g_{-3}=b(\mathfrak g_{-1},\mu)_{-3}/\mathfrak a_{-3}$
is isomorphic to
\[
(\mathfrak e\otimes \mathfrak e\otimes \mathfrak f)/
\Lambda^2(\mathfrak e)\otimes\mathfrak f\oplus
(\mathfrak e\otimes \mathfrak f\otimes \mathfrak f)/
\mathfrak e\otimes\Lambda^2(\mathfrak f)
\cong
S^2(\mathfrak e)\otimes
\mathfrak f\oplus S^2(\mathfrak f)\otimes \mathfrak e.
\]
This completes the proof.
\end{proof}

\begin{proposition}\label{prop8.4}
Let $\mathfrak m=\bigoplus\limits_{p<0}\mathfrak g_p$ be a
pseudo-product FGLA of the $\mu$-th kind with
pseudo-product structure $(\mathfrak e,\mathfrak f)$, where $\mu\geqq2$.
We denote by $\mathfrak c$
the centralizer of $\mathfrak g_{-2}$ in $\mathfrak g_{-1}$.
Let $\gla g$ be the prolongation of $(\mathfrak m;\mathfrak e,\mathfrak f)$.
Assume that $\mathfrak g_0$ is isomorphic to $\gl(\mathfrak e)\oplus\gl(\mathfrak f)$ as a Lie algebra.
\begin{enumerate}\itemsep=0pt
\item If $\mu=2$, then $\mathfrak m=\bigoplus\limits_{p<0}\mathfrak g_p$
be a free
pseudo-product FGLA.
\item If $\mu\geqq3$ and $\mathfrak c\ne\{0\}$,
then $(\mathfrak m;\mathfrak e,\mathfrak f)$ is not a free pseudo-product
FGLA.
\item If $\mu=3$ and $\mathfrak c=\{0\}$,
then $(\mathfrak m;\mathfrak e,\mathfrak f)$ is a free pseudo-product
FGLA.
\end{enumerate}
\end{proposition}

\begin{proof}
Let
$\check{\mathfrak m}=\bigoplus_{p<0}\limits\check{\mathfrak g}_p$
be the free pseudo-product FGLA of type $(m,n,\mu)$
with pseudo-product structure
$(\check{\mathfrak e},\check{\mathfrak f})$ such that
$\check{\mathfrak g}_{-1}=\mathfrak g_{-1}$,
$\check{\mathfrak e}=\mathfrak e$ and $\check{\mathfrak f}=\mathfrak f$.
Let $\check{\mathfrak g}=\bigoplus_{p\in\mathbb Z}\limits\check{\mathfrak g}_p$  be the prolongation of $(\check{\mathfrak m};\check{\mathfrak e},\check{\mathfrak f})$.
There exists a GLA epimorphism $\varphi$
of $\check{\mathfrak m}$ onto $\mathfrak m$ such that the restriction
$\varphi|\check{\mathfrak g}_{-1}$ is the identity mapping.
Since the mapping
$\check{\mathfrak g}_0\ni D\mapsto
(D|\mathfrak e,D|\mathfrak f)\in\gl(\mathfrak e)\times\gl(\mathfrak f)$
is an isomorphism,
$\varphi$ can be extended to be a homomorphism $\check{\varphi}$
of $\bigoplus\limits_{p\leqq0}\check{\mathfrak g}_{p}$ onto
$\bigoplus\limits_{p\leqq0}\mathfrak g_p$.
Let $\mathfrak a$ be the kernel of $\check{\varphi}$;
then $\mathfrak a$ is a graded ideal of
$\bigoplus\limits_{p\leqq0}\check{\mathfrak g}_{p}$.
We set $\mathfrak a_p
=\mathfrak a\cap \check{\mathfrak g}_{p}$; then
$\mathfrak a=\bigoplus\limits_{p\leqq 0}\mathfrak a_p$.
Since the restriction of $\check{\varphi}$ to
$\check{\mathfrak g}_{-1}\oplus \check{\mathfrak g}_0$ is injective,
$\mathfrak a_p=\{0\}$ for $p\geqq-1$. Also
each $\mathfrak a_p$ is a $\check{\mathfrak g}_0$-submodule of
$\check{\mathfrak g}_p$.
Since the $\check{\mathfrak g}_0$-module
$\check{\mathfrak g}_{-2}$ is irreducible (Proposition~\ref{prop8.3}~(2)),
$\varphi|\mathfrak g_{-2}$ is injective.
If $\mu=2$, then $\varphi$ is an isomorphism. This proves the assertion~(1).
Now we assume that $\mu\geqq3$. Then
\[
\check{\mathfrak g}_{-3}=[[\mathfrak e,\mathfrak f],\mathfrak f]\oplus
[[\mathfrak e,\mathfrak f],\mathfrak e].
\]
Since $\check{\mathfrak g}_0$-modules
$[[\mathfrak e,\mathfrak f],\mathfrak f]$ and
$[[\mathfrak e,\mathfrak f],\mathfrak e]$ are irreducible and not isomorphic
to each other (Proposition \ref{prop8.3} (3)),
one of the following cases occurs:
$(i)$~$\mathfrak a_{-3}=[[\mathfrak e,\mathfrak f],\mathfrak f]$;
$(ii)$~$\mathfrak a_{-3}=[[\mathfrak e,\mathfrak f],\mathfrak e]$;
$(iii)$~$\mathfrak a_{-3}=\{0\}$.
If $\mathfrak a_{-3}=[[\mathfrak e,\mathfrak f],\mathfrak f]$ (resp.\
$\mathfrak a_{-3}=[[\mathfrak e,\mathfrak f],\mathfrak e]$), then
$\mathfrak c=\mathfrak f$
(resp. $\mathfrak c=\mathfrak e$).
Also
since $\mathfrak g_0$-modules
$\mathfrak e$,
$\mathfrak f$ are irreducible and not isomorphic to each other,
one of the following cases occurs:
$(i)$~$\mathfrak c=\mathfrak e$;
$(ii)$~$\mathfrak c=\mathfrak f$;
$(iii)$~$\mathfrak c=\{0\}$.
If $\mathfrak c=\mathfrak e$
(resp. $\mathfrak c=\mathfrak f$),
then $\mathfrak a_{-3}=[[\mathfrak e,\mathfrak f],\mathfrak e]$
(resp. $\mathfrak a_{-3}=[[\mathfrak e,\mathfrak f],\mathfrak f]$).
In this case, $\varphi$ is not injective.
Hence $(\mathfrak m;\mathfrak e,\mathfrak f)$ is not free.
If $\mathfrak c=\{0\}$,
then $\mathfrak a_{-3}=\{0\}$.
Hence $\varphi|\check{\mathfrak g}_{-3}$ is an isomorphism.
In particular, if $\mu=3$, then
$(\mathfrak m;\mathfrak e,\mathfrak f)$ is free.
\end{proof}

\begin{example}
Let $V$ and $W$ be f\/inite-dimensional vector spaces and $k\geqq1$.
We set
\begin{gather*}
\mathfrak C^{k}(V,W) =\bigoplus\limits_{p=-k-1}^{-1}\mathfrak C^{k}(V,W)_p, \\
\mathfrak C^{k}(V,W)_p=W\otimes S^{k+p+1}(V^*), \qquad
 -k-1\leqq p\leqq -2, \\
\mathfrak C^{k}(V,W)_{-1} =V\oplus (W\otimes S^{k}(V^*)).
\end{gather*}
The bracket operation of $\mathfrak C^{k}(V,W)$ is def\/ined as follows:
\begin{gather*}
  [W,V]=\{0\},\qquad [V,V]=\{0\}, \qquad [W\otimes S^{r}(V^*),W\otimes S^{s}(V^*)]=\{0\}, \\
[w\otimes s_r,v]=w\otimes (v\lrcorner\,s_r)\qquad \text{for}\  v\in V,\ w\in W,
\ s_r\in S^r(V^*).
\end{gather*}
Equipped with this bracket operation, $\mathfrak C^{k}(V,W)$ becomes
a pseudo-product FGLA of the $(k+1)$-th kind
with pseudo-product structure $(V,W\otimes S^{k}(V^*))$,
which is called
{\it the contact algebra of order $k$ of bidegree} $(n,m)$, where
$n=\dim V$ and $m=\dim W$ (cf.~\cite[p.~133]{yam82:1}).
We assume that $\mathfrak C^{k}(V,W)$ is a free pseudo-product FGLA.
Since
\begin{gather*}
\dim \mathfrak C^{k}(V,W)_{-2}=m\binom{n+k-2}{k-1}, \qquad
\dim V\dim (W\otimes S^{k}(V^*))= nm\binom{n+k-1}{k},
\end{gather*}
we get $n=1$.
Since $W\otimes S^{k}(V^*)$ is contained in the centralizer of
$\mathfrak C^{k}(V,W)_{-2}$
in $\mathfrak C^{k}(V,W)_{-1}$, we get $k=1$.
Thus we obtain that $\mathfrak C^{k}(V,W)$ is a free pseudo-product FGLA
if and only if $k=1$, $n=1$.
\end{example}

\begin{example}\label{ex8.2}
Let $\gla g$ be a f\/inite-dimensional simple GLA of type
$(A_{m+n},\{\alpha_m,\alpha_{m+1}\})$.
We set $\mathfrak e=\mathfrak g_{-1}^{(m)}$, $\mathfrak f=\mathfrak g_{-1}^{(m+1)}$. Then $(\mathfrak g_-;\mathfrak e,\mathfrak f)$
is a pseudo-product FGLA.
Since $\dim \mathfrak e=m$, $\dim\mathfrak f=n$ and
$\dim\mathfrak g_{-2}=mn$, the pseudo-product FGLA
$(\mathfrak g_-;\mathfrak e,\mathfrak f)$ is a free pseudo-product FGLA
of type $(m,n,2)$ (Proposition 8.3 (2)). Also
$\gla g$ is the prolongation of $\mathfrak g_-$ except for the following
cases (see \cite{yam93:1}):
\begin{enumerate}\itemsep=0pt
\item $m=n=1$. In this case,
the prolongation of $\mathfrak g_-$ is isomorphic to $K(1)$.
\item $m=1$ or $n=1$ and $l=\max\{m,n\}\geqq2$. In this case,
the prolongation of $\mathfrak g_-$ is isomorphic to $W(l+1;\bm s)$,
where $\bm s=(1,2,\dots,2)$.
\end{enumerate}
\end{example}

\begin{example}
Let $V$ and $W$ be f\/inite-dimensional vector spaces
such that $\dim V=m\geqq1$ and $\dim W=n\geqq1$.
We set
\begin{gather*}
\mathfrak g_{-1}  =V\oplus W, \qquad \mathfrak g_{-2}=V\otimes W, \\
\mathfrak g_{-3}  = V\otimes S^2(W)\oplus S^2(V)\otimes W,
\qquad \mathfrak m=\mathfrak g_{-1}\oplus \mathfrak g_{-2}\oplus
\mathfrak g_{-3}.
\end{gather*}
The bracket operation of $\mathfrak m$ is def\/ined as follows:
 \begin{gather*}
  [\mathfrak g_{-3},\mathfrak g_{-1}\oplus \mathfrak g_{-2}]
=[\mathfrak g_{-2},\mathfrak g_{-2}]=\{0\},  \qquad [V,V]=[W,W]=\{0\},\\
[v,w]=-[w,v]=v\otimes w, \qquad
  [v,v'\otimes w]=-[v'\otimes w,v]=v\circledcirc v'\otimes w,
\\ [v\otimes w,w']=-[w',v\otimes w]=v\otimes w\circledcirc w',
\end{gather*}
where $v,v'\in V$ and $w,w'\in W$.
Equipped with this bracket operation, $\mathfrak m$ becomes
a free pseudo-product FGLA of type $(m,n,3)$ with pseudo-product structure
$(V,W)$ (Proposition~\ref{prop8.3}).
\end{example}
